\makeabstract
    {
    Auditory Processing and Attention to Speech in Minimally Verbal Autism: A Pilot Study
    }
    {
    Evie Brimberg\textsuperscript{1}, Theo Vanneau\textsuperscript{2}, Chloe Brittenham\textsuperscript{2}, Jackson Mitchell\textsuperscript{2}, Sydney Kornfeld\textsuperscript{2}, Elizabeth Akinyemi\textsuperscript{2}, Sophie Molholm\textsuperscript{2}
    }
    {
    \textsuperscript{1} Amherst College, Amherst, MA\\
\textsuperscript{2} Departments of Pediatrics and Neuroscience, Albert Einstein College of Medicine, Bronx, NY
    }
    {
    Autism spectrum disorder (ASD) involves substantial heterogeneity in language ability and cognitive function. However, little is known about minimally verbal autism (MVA), which comprises an estimated 30\% of the autistic population, as most research has focused on individuals who communicate through spoken language and demonstrate average to above-average intelligence. This pilot study examines how auditory processing and attentional orienting to novel auditory stimuli relate to communication and cognitive ability across minimally verbal autism (MVA), verbal autism (VA), and typically developing (TD) groups, aiming to uncover neurophysiological factors contributing to MVA.
    }
    {
    Scalp EEG was recorded in 18 children ages 8{\textendash}14 (MVA, n=6; VA, n=7; TD, n=5) during a passive auditory oddball paradigm presenting task-irrelevant tones and consonant-vowel (CV) sounds while participants watched a muted video. Standard and deviant stimuli belonged to the same stimulus class per block; novel stimuli were drawn from a different class. Participants also completed measures of receptive (PPVT) and expressive (EVT) language, cognitive ability (IQ), and adaptive behavior (Vineland). Mismatch negativity (MMN) and P3a components were derived from novel-minus-standard difference waves.
    }
    {
    Difference waves showed an early fronto-central deflection consistent with a MMN ($\sim$175 ms for tones, $\sim$250 ms for CV sounds, though this classification is tentative for CV), followed by a P3a ($\sim$300 ms tone, $\sim$340 ms CV). MVA participants showed a reduced P3a response to novel CV sounds relative to TD (p=.013, Holm-corrected), consistent with reduced automatic attentional orienting to unexpected speech sounds. Neural component amplitude showed preliminary associations with clinical measures of communication: tone MMN amplitude correlated negatively with Vineland-3 Communication scores (r={\textendash}.55, p=.022), and CV P3a amplitude correlated positively with Vineland-3 Communication scores (r=.61, p=.009). 
    }
    {
    These preliminary findings suggest that reduced attentional orienting to novel speech sounds, indexed by P3a, relates to communication ability across the autism spectrum. Data collection is ongoing to increase power for group comparisons and correlational analyses, and future analyses will extend to word-listening and continuous naturalistic-speech conditions to examine auditory processing under more naturalistic listening settings.
    }

    \newpage
    